\chapter{Nghiên cứu liên quan}
\label{ch:lienquan}

\section{Tổng quan các kỹ thuật bảo vệ quỹ đạo}

Chow và Mokbel~\cite{chow2011trajectory} hệ thống hóa các kỹ thuật bảo vệ quỹ đạo
trong LBS liên tục và xuất bản dữ liệu: spatial cloaking, mix-zones, path
confusion và dummy. Hạn chế chung của nhóm kỹ thuật này là không có đảm bảo hình
thức và dễ vỡ trước tấn công dựa trên tri thức nền. Hai công trình của nhóm
nghiên cứu HCMUT --- PP$^{ST}$-tree~\cite{phan2012ppst} và
KUR-Algorithm~\cite{phan2015kur} --- bảo vệ quỹ đạo ở tầng cơ sở dữ liệu bằng
k-anonymity trên các vùng dự đoán chuyển động, đặt nền cho hướng tiếp cận nhận
biết ngữ cảnh đường bộ mà luận văn kế thừa.

\section{Geo-I và các mở rộng}

Geo-I~\cite{andres2013geo} cung cấp đảm bảo định lượng per-point nhưng có ba hạn
chế đã được chỉ rõ trong văn liệu: không gian được đối xử đồng nhất
(khắc phục bởi elastic metrics~\cite{chatzikokolakis2015elastic}); metric Euclid
không khớp ngữ cảnh đường bộ; và đầu ra có thể rơi vào vị trí phi thực tế.
Bordenabe và cộng sự~\cite{bordenabe2014optimal} xây dựng cơ chế Geo-I tối ưu
tiện ích bằng quy hoạch tuyến tính. Oya và cộng sự~\cite{oya2017back} chỉ ra
Geo-I đơn lẻ có thể để lọt rủi ro suy luận Bayes cao, khuyến nghị bổ sung độ đo
sai số suy luận kỳ vọng.

\paragraph{Geo-I trên lưới đường.} Gần với luận văn nhất,
Takagi và cộng sự~\cite{takagi2019geo} đề xuất
\emph{Geo-Graph-Indistinguishability} (GG-I): miền đầu ra là tập đỉnh đồ thị
đường, metric là khoảng cách đường đi ngắn nhất $d_s$, và cơ chế
Graph-Exponential Mechanism $\Pr[o \mid v] \propto \exp(-\tfrac{\eps}{2} d_s(v,o))$.
Họ chứng minh mọi cơ chế $\eps$-Geo-I hậu xử lý bằng phép chiếu về đỉnh gần nhất
cũng thỏa $\eps$-GG-I, và GEM trội hơn Laplace-rồi-snap trên trade-off riêng
tư--tiện ích. Cơ chế REM của luận văn (Mục~\ref{sec:rem}) là phiên bản dùng
metric Euclid của hướng này, đổi độ chính xác metric lấy chi phí tính toán
$O(|V|)$ thay vì Dijkstra trên toàn đồ thị.

\section{Tương quan thời gian}

Xiao và Xiong~\cite{xiao2015protecting} chỉ ra Geo-I per-point suy giảm dưới
tương quan thời gian: kẻ tấn công lan truyền prior qua ma trận chuyển Markov và
lọc nhiễu như một bộ lọc Bayes. Họ đề xuất $\delta$-location set và cơ chế PIM.
PIVE~\cite{yu2017pive} kết hợp DP với ràng buộc sai số suy luận qua Protection
Location Set hai pha; tuy nhiên Zhang và cộng sự~\cite{zhang2021pive} đã chứng
minh cấu trúc PLS thích nghi của PIVE không đạt được các đảm bảo như tuyên bố.
PTPPM~\cite{cao2024ptppm} thay cơ chế lũy thừa bằng Permute-and-Flip trên PLS.
Bản preprint gần nhất của nhóm này~\cite{min2025road} (11/2025, chưa qua bình
duyệt) kết hợp lưới đường với tương quan thời gian --- xác nhận đây là khoảng
trống nghiên cứu đang mở mà luận văn nhắm tới, với cách tiếp cận đơn giản hơn
và đảm bảo sạch hơn (trọng số công khai, Định lý~\ref{thm:trem}).

\section{Học sâu và tấn công tái tạo}

Các mô hình sinh quỹ đạo (LSTM-TrajGAN~\cite{rao2020lstm},
DiffTraj~\cite{zhu2023difftraj}, Diff-RNTraj trên lưới đường) cho quỹ đạo giả
trông thực nhưng không có đảm bảo DP; huấn luyện kèm DP-SGD làm tiện ích suy giảm
mạnh~\cite{buchholz2024sok}. Quan trọng với luận văn là tấn công
RAoPT~\cite{buchholz2022raopt}: mạng LSTM học ``chữ ký'' của nhiễu DP (điểm rơi
ngoài đường, trên mặt nước) và tái tạo quỹ đạo gốc, giảm hơn 65\% khoảng cách
Euclid/Hausdorff về quỹ đạo thật trên T-Drive với $\eps \le 1$. Đây là luận cứ
thực nghiệm trung tâm cho thiết kế ``đầu ra trên đường theo cấu trúc'' của REM.

\section{Dataset và độ đo chuẩn}
\label{sec:datasetmetrics}

Khảo sát các công trình 2013--2026 cho thấy nhánh trajectory-level dùng áp đảo
\textbf{GeoLife}~\cite{zheng2010geolife} (182 người dùng, 17.621 quỹ đạo GPS,
Bắc Kinh, 91\% mật độ 1--5 giây) --- đây cũng là dataset của PIVE, PIM, RAoPT.
Bộ độ đo chuẩn mà hội đồng/bình duyệt kỳ vọng gồm: \emph{(privacy)} sai số suy
luận kỳ vọng dưới tấn công Bayes tối ưu~\cite{shokri2011quantifying} và tấn công
theo dõi Markov/HMM~\cite{xiao2015protecting}; \emph{(utility)} quality loss kỳ
vọng và độ chính xác truy vấn k-NN POI; \emph{(realism)} tỷ lệ điểm trên đường và
tính hợp lý tốc độ giữa hai điểm công bố liên tiếp. Chương~\ref{ch:thucnghiem}
hiện thực đầy đủ bộ độ đo này.
